\documentclass[../../../main.tex]{subfiles}
\begin{document}
Solids that can be described as crystal are called polycrystalline; these material are actually composed of small crystalline grains.
This collection of crystal--or grains--have their own orientation.
The interface of two grains with different orientation are called grain boundary and this can often make the crystal have more mechanical strength.

On the other hand, solids that  characterized by the absence of any long-range order are called amorphous.
Without order, their atomic structure cannot form grain, thus they do not have grain boundary.

\subsection{Mathematical Description}
\subsubsection{Primitive vector.}
Primitive vector, also called primitive translation vector, is one of the fundamental displacement vectors that generate the Bravais lattice by translation.
We arrange the three primitive vectors $\mathbf{a }_1,\mathbf{a }_2,\mathbf{a }_3$ as columns of a $3 \times 3$ matrix
\begin{equation*}
    A=\begin{bmatrix}
        \mathbf{a }_1 & \mathbf{a }_2 & \mathbf{a }_3
    \end{bmatrix}
\end{equation*}
where each vector is understood as
\begin{equation*}
    \mathbf{a}_1=
    \begin{bmatrix}
        a_{11} \\
        a_{12} \\
        a_{13} \\
    \end{bmatrix}\qquad
    \mathbf{a}_2=
    \begin{bmatrix}
        a_{21} \\
        a_{22} \\
        a_{23} \\
    \end{bmatrix}\qquad
    \mathbf{a}_3=
    \begin{bmatrix}
        a_{31} \\
        a_{32} \\
        a_{33} \\
    \end{bmatrix}
\end{equation*}
such that
\begin{equation*}
    A=\begin{bmatrix}
        a_{11} & a_{12} & a_{13} \\
        a_{21} & a_{22} & a_{23} \\
        a_{31} & a_{32} & a_{33} \\
    \end{bmatrix}
\end{equation*}

\subsubsection{Bravais lattice.}
Defined as an infinite array of discrete points with identical surroundings.
These discrete points are generated by the primitive generating vector
\begin{equation*}
    \mathbf{R} =m \mathbf{a }_1+n \mathbf{a }_2 +o \mathbf{a }_3
\end{equation*}
If we were to consider this definition only, we will have the definition of mathematical lattice.
The Bravais lattice then put another restriction: any lattice point can be chosen as the origin, and the set of relative positions of all other lattice points remains exactly the same.
The equivalence of all point, in other words.

In the case of simple 2D lattice, the magnitude of primitive vector $\mathbf{a }_n$ is the lattice constant.
However, for 3D lattice, the primitive vector $\mathbf{a }_n$ do not align with the cube edges--for the case of FCC and BCC.
In this case, the lattice constant is defined as the length of the cube edges.

\subsubsection{Lattice plane.}
A lattice plane in a Bravais lattice is the infinite set of lattice points that lie on a plane spanned by two linearly independent lattice vectors.

\subsubsection{Unit cells.}
\begin{figure*}[b]
    \centering
    \normfig{../../../Rss/QM/ManyBody/unitcells.jpg}
    \caption*{Figure: Primitive and nonprimitive unit cells}
\end{figure*}
Defined as any volume of space when translated through all the vector in Bravais lattice, fills space without overlap and without leaving voids.
Primitive unit cells contain only one lattice point, whereas nonprimitive unit cells contain more than  one unit cells.
A special choice of the primitive unit cell is the
Wigner-Seitz cell that is defined as the closest region of space to one given lattice point than to any other.

The volume of the unit cells can be determined by the scalar triple product rule
\begin{equation*}
    V=\mathbf{a}_1 \cdot \mathbf{a }_2 \times \mathbf{a }_3
\end{equation*}

\begin{table}[h]
    \centering
    \caption*{Table: Difference of primitive and nonprimitive unit cells}
    \begin{tabular}{ cc>{\centering\arraybackslash}p{2cm}}
        \toprule
                       & Primitive           & Nonprimitive                                                                                   \\
        \midrule
        Volume         & Smaller             & Bigger                                                                                         \\
        Lattice points & One                 & More than one                                                                                  \\
        Function       & Standard definition & Structure representation--such as FCC, HCP, and BCC--whose primitive unit cells can be derived \\
        \bottomrule
    \end{tabular}
\end{table}

\subsubsection{Basis.}
Defined as a set of one or more atoms associated with each lattice point.
Using both Bravais lattice and basis, the crystal structure can be described.
The position of all atom--also called crystal points--is written as
\begin{equation*}
    \text{Crystal points}=\mathbf{R}+\mathbf{r}_i
\end{equation*}
with $\mathbf{r}_i$ as the position of the $i$-th atom in the basis relative to that lattice points.
\begin{figure*}[b]
    \centering
    \normfig{../../../Rss/QM/ManyBody/basis.jpeg}
    \caption*{Figure: Basis}
\end{figure*}

\subsubsection{APF.}
The Atomic Packing Factor (APF) quantifies how efficiently atoms are packed in a crystal structure
\begin{equation*}
    \text{APF}=\frac{\text{Effective volume}}{\text{Total volume of unit cell}}
\end{equation*}
The effective volume also depend on effective volume, that is $V_\text{eff}=N \cdot V_\text{atom}$, with $N$ as the number of effective atom.

\subsubsection{Miller indices.}
Miller indices $(hkl)$ is used to characterize lattice planes.
In geometric terms, the integers $(h,k,l,)$ are related to the intercepts of the plane with the real-space primitive vectors $(\mathbf{a }_1,\mathbf{a }_2,\mathbf{a }_3)$ by
\begin{equation*}
    (h,k,l)=\left( \frac{1 }{p},\frac{1 }{q}, \frac{1 }{r} \right)
\end{equation*}
where $p\mathbf{a }_1,q\mathbf{a }_2,r\mathbf{a }_3$ is the plane intercepts of the axes.
Also, bar is conventionally used when denoting negative Miller indices.

\subsubsection{Reciprocal lattice.}
For a three-dimensional Bravais lattice with primitive vectors $\mathbf{a}_1, \mathbf{a}_2, \mathbf{a}_3$, the reciprocal lattice primitive vectors are
\begin{equation*}
    \mathbf{b}_1 = 2 \pi \frac{\mathbf{a}_2 \times \mathbf{a}_3}{\mathbf{a}_1 \cdot (\mathbf{a}_2 \times \mathbf{a}_3)}, \quad
    \mathbf{b}_2 = 2 \pi \frac{\mathbf{a}_3 \times \mathbf{a}_1}{\mathbf{a}_1 \cdot (\mathbf{a}_2 \times \mathbf{a}_3)}, \quad
    \mathbf{b}_3 = 2 \pi \frac{\mathbf{a}_1 \times \mathbf{a}_2}{\mathbf{a}_1 \cdot (\mathbf{a}_2 \times \mathbf{a}_3)}
\end{equation*}
The reciprocal lattice vector can be expressed as
\begin{equation*}
    \mathbf{G} = h \mathbf{b}_1 + k \mathbf{b}_2 + l \mathbf{b}_3, \quad h, k, l \in \mathbb{Z}
\end{equation*}
The reciprocal lattice is fundamental in diffraction phenomena, where the condition for constructive interference is
\begin{equation*}
    \Delta \mathbf{k} = \mathbf{G}.
\end{equation*}
For given Bravais lattice, we can also write
\begin{equation*}
    e^{i \mathbf{G} \cdot \mathbf{R}} = 1, \qquad
    \mathbf{G} \cdot \mathbf{R} = 2 \pi n
\end{equation*}

\subsection{Lattice Example}
\subsubsection{2 dimensional.}
\begin{figure*}[h]
    \centering
    \fullfig{../../../Rss/QM/ManyBody/2DBravais-3.png}
    \caption*{Figure: 2D Bravais lattice. Respectively square, rectangular, oblique, centered rectangular, and hexagonal}
\end{figure*}

There are five distinct types of two-dimensional Bravais lattices.
They are classified by their lattice vectors and symmetries
\begin{enumerate}
    \item \textbf{Square.} Lattice vectors with equal lengths and orthogonal angle $\gamma =90^\circ$.
          This structure can be generated by
          \begin{equation*}
              A_{\text{square}} = a
              \begin{bmatrix}
                  1 & 0 \\
                  0 & 1
              \end{bmatrix}
          \end{equation*}
    \item \textbf{Rectangular.} Lattice vectors with unequal lengths and orthogonal angle $\gamma =90^\circ$.
          This structure can be generated by
          \begin{equation*}
              A_{\text{rect}} =
              \begin{bmatrix}
                  a & 0 \\
                  0 & b
              \end{bmatrix}
          \end{equation*}
    \item \textbf{Oblique.} Lattice vectors with unequal lengths and arbitrary angle $\gamma \neq 0$.
          This structure can be generated by
          \begin{equation*}
              A_{\text{oblique}} =
              \begin{bmatrix}
                  a & b\cos\gamma \\
                  0 & b\sin\gamma
              \end{bmatrix}
          \end{equation*}
    \item \textbf{Centered rectangular.} Same as rectangular, but with an additional lattice point at the center of each rectangle.
          This structure can be generated by
          \begin{equation*}
              A_{\text{c-rect}} = \frac{1}{2}
              \begin{bmatrix}
                  a & a  \\
                  b & -b
              \end{bmatrix}
          \end{equation*}
    \item \textbf{Hexagonal.} Lattice vectors with equal lengths and orthogonal angle $\gamma =120^\circ$.
          This structure can be generated by
          \begin{equation*}
              A_{\text{hex}} = a
              \begin{bmatrix}
                  1 & 1/2          \\
                  0 & \sqrt{3}/{2}
              \end{bmatrix}
          \end{equation*}
\end{enumerate}

\subsubsection{3 dimensional.}
There are 7 classification for crystal system, which totaling into 14 distinct Bravais lattice.
They are classified based on the lengths of the lattice parameters $a,b,c$; interaxial angles $\alpha,\beta,\gamma$; and possible centering types: primitive $P$, body-centered $I$, face-centered $F$, and base-centered $C$.
\begin{table}[h]
    \centering
    \caption*{Table: The seven different basis-vector systems or crystal systems}
    \begin{tabular}{cccc}
        \toprule
        Basis vector      & Angle                                              & Crystal system & Type      \\
        \midrule
        $a \neq b \neq c$ & $\alpha \neq \beta\neq \gamma \neq 90 ^\circ$      & Triclinic      & $P$       \\
        $a \neq b \neq c$ & $\alpha =\gamma =90 ^\circ,\; \beta\neq 90 ^\circ$ & Monoclinic     & $P,C$     \\
        $a \neq b \neq c$ & $\alpha = \beta= \gamma 90 ^\circ$                 & Orthorombic    & $P,C,I,F$ \\
        $a =b \neq c$     & $\alpha = \beta= \gamma=90 ^\circ$                 & Tetragonal     & $P,I$     \\
        $a =b \neq c$     & $\alpha = \beta= 90 ^\circ,\;\gamma=120 ^\circ$    & Hexagonal      & $P$       \\
        $a =b = c$        & $\alpha=\beta=\gamma \neq 90 ^\circ$               & Rhombohedral   & $R$       \\
        $a =b = c$        & $\alpha=\beta =\gamma90 ^\circ$                    & Cubic          & $P,I,F$   \\
        \bottomrule
    \end{tabular}
\end{table}

\begin{figure*}[]
    \centering
    \fullfig{../../../Rss/QM/ManyBody/3DBravais.jpg}
    \caption*{Figure: The 14 three-dimensional Bravais lattices}
\end{figure*}

\subsection{Crystal Structure}
\begin{figure*}[]
    \centering
    \normfig{../../../Rss/QM/ManyBody/CubicStructure.jpg}
    \caption*{Figure: Three type of cubic structure: simple cubic, body-cetered cubic, and face centered cubic}
\end{figure*}

\subsubsection{Simple cubic.}
This is a cubic atomic structure where atom are positioned at the eight corners of the cube and each is shared among eight neighboring unit cells.
The number of effective atomic number is
\begin{equation*}
    \frac{1 }{8 }\cdot 8=1
\end{equation*}
The packing density then is calculated as follows
\begin{equation*}
    \text{APF}=\frac{1 \cdot 4 \pi r^3/3 }{(2r)^3}\approx 0.52
\end{equation*}
where we construct the unit cells in terms of the atom radius.

The primitive vector is written as
\begin{equation*}
    A_{\text{SC}} = a
    \begin{bmatrix}
        1 & 0 & 0 \\
        0 & 1 & 0 \\
        0 & 0 & 1
    \end{bmatrix}
\end{equation*}

\subsubsection{Body-centered cubic.}
This is a cubic atomic structure where atom are positioned at the eight cube corners and one atom at the cube center.
The number of effective atomic number is
\begin{equation*}
    \frac{1 }{8 }\cdot 8+1=2
\end{equation*}
The packing density then is calculated as follows
\begin{equation*}
    \text{APF}=\frac{2 \cdot 4\pi r^3/2 }{(4r/\sqrt{3})^2}\approx 0.68
\end{equation*}
where we construct the unit cells in terms of the atom radius
\begin{equation*}
    \sqrt{3} a_0=4r \implies a_0=\frac{4r }{\sqrt{3}}
\end{equation*}

The primitive vector is written as
\begin{equation*}
    A_{\text{BCC}} = \frac{a}{2}
    \begin{bmatrix}
        -1 & 1  & 1  \\
        1  & -1 & 1  \\
        1  & 1  & -1
    \end{bmatrix}
\end{equation*}

\subsubsection{Face-centered cubic.}
This is a cubic atomic structure where atom are positioned at the eight cube corners and at the centers of all six cube faces.
Each face-centered atom is shared between two unit cells.
The number of effective atomic number is
\begin{equation*}
    \frac{1 }{8 }\cdot 8+\frac{1 }{2 }\cdot 6=4
\end{equation*}
The packing density then is calculated as follows
\begin{equation*}
    \text{APF}=\frac{4 \cdot 4\pi r^3/\sqrt{2} }{(4r/\sqrt{2})^3}\approx 0.74
\end{equation*}
where we construct the unit cells in terms of the atom radius.
\begin{equation*}
    \sqrt{2} a_0=4r \implies a_0=\frac{4r }{\sqrt{2}}
\end{equation*}

The primitive vector is written as
\begin{equation*}
    A_{\text{FCC}} = \frac{a}{2}
    \begin{bmatrix}
        0 & 1 & 1 \\
        1 & 0 & 1 \\
        1 & 1 & 0
    \end{bmatrix}
\end{equation*}

\begin{figure*}[]
    \centering
    \fullfig{../../../Rss/QM/ManyBody/CubicAPF.jpg}
    \caption*{Figure: The relationships between the atomic radius and the lattice parameter in cubic systems}
\end{figure*}

Along with HCP, FCC is the densest possible arrangements, that is both has the APF of 0.74.
This is also why FCC and HCP are called Close-Packed Structures.
To construct FCC, one start with hexagonal structure, the densest packing; this is the $A$ layer.
Then stack another hexagonal layer on the depression of the first, this is the $B$ layer.
The third layer, then, consist of hexagonal layer that occupy depressions not aligned with $A$ or $B$; this layer is $C$.
Thus, FCC has the sequence of $ABC\;ABC\dots$.

\subsubsection{Hexagonal close packed.}
The Hexagonal Close-Packed (HCP) structure is one of the two possible three-dimensional close-packed crystal structures (the other being Face-Centered Cubic, FCC).
The number of effective atomic number is
\begin{equation*}
    \frac{1 }{6 }\cdot 12+\frac{1 }{2 }\cdot 2+1 \cdot 3=6
\end{equation*}
where the corner atom being shared by six unit cells, and the face atom being shared by two unit cells.
There are also three atoms inside the unit cell.
To calculate the APF, first consider the volume of the hexagonal base
\begin{equation*}
    A=\frac{3 \sqrt{3 }}{2 }a^2
\end{equation*}
with $a=2r$ as the hexagonal side length.
For a unit cell with height $c$, the volume is then
\begin{equation*}
    V = \left|\mathbf{a}_1 \cdot (\mathbf{a}_2 \times \mathbf{a}_3)\right|=
    \left|
    \det
    \begin{pmatrix}
        a             & 0                    & 0 \\[4pt]
        -\tfrac{a}{2} & \tfrac{\sqrt{3}a}{2} & 0 \\[4pt]
        0             & 0                    & c
    \end{pmatrix}
    \right|
    = \frac{\sqrt{3}}{2}\,a^{2}c.
\end{equation*}
For ideal close packing,
\begin{equation*}
    \frac{c }{a }=\sqrt{\frac{8 }{3 }}\approx1.633
\end{equation*}
or in terms of radius $r$
\begin{equation*}
    c=3.266r
\end{equation*}
By this, we can express the volume of the unit cells as
\begin{equation*}
    V\approx33.94r^3
\end{equation*}
and the packing density as
\begin{equation*}
    \text{APF}=\frac{6 \cdot 4\pi r^3/3 }{33.94r^3}\approx 0.74
\end{equation*}

The structure can be generated by
\begin{equation*}
    A_{\text{HCP}} =
    \begin{bmatrix}
        a & -a/2          & 0 \\
        0 & \sqrt{3}/{2}a & 0 \\
        0 & 0             & c
    \end{bmatrix}
\end{equation*}

The method to construct HCP is similar with FCC: start with hexagonal layer $A$ and stack another hexagonal layer $B$ on the depressions.
Unlike HCP, the third layer consist of hexagonal layer aligned with the first layer.
Thus, HCP has the sequence of $AB\;AB\dots$.

\begin{figure*}[]
    \centering
    \normfig{../../../Rss/QM/ManyBody/ClosedPackingStr.png}
    \caption*{Figure: Close packing of spheres leading to the HCP $(AB\;AB\dots)$ and fcc $(ABC\;ABC\dots)$ structures}
\end{figure*}

\subsection{Covalent Structure}
In covalent structures, the atoms’ valence electrons are not completely delocalized but shared between neighboring atoms and the bond length and direction are far more important than the packing density.
\subsubsection{Graphene.}
Graphene is a covalent structure with single sheet of carbon atoms in a honeycomb lattice structure.
The carbon atoms in graphene are connected by sp$^2$ hybrid bonds, enclosing an angle of 120$^\circ$.

\subsubsection{Graphite.}
The parent material of graphene is graphite, a stack of graphene sheets that are weakly bonded to each other.
Carbon atoms are sp$^2$ hybridized, forming hexagonal sheets with strong covalent bonding within each layer and weak van der Waals interactions between layers.

\subsubsection{Diamond.}
In diamond, the carbon atoms form sp$^3$ -type bonds and each atom has four nearest neighbors in a tetrahedral configuration.
Interestingly, the diamond structure can also be described as an FCC Bravais lattice with a basis of two atoms.
\begin{figure*}[]
    \centering
    \normfigL{../../../Rss/QM/ManyBody/Covalent}
    \caption*{Figure: Structures for (a) graphene, (b) graphite, and (c) diamond.}
\end{figure*}

\subsection{Ionic Structure}
Ionic materials must have crystal structures that ensure electrical neutrality, yet permit ions of different sizes to be packed efficiently.

An anion is an atom or group of atoms that has gained one or more electrons, giving it a net negative charge; in other hand, cation is an atom or group of atoms that has lost one or more electrons, giving it a net positive charge.
Their electrostatic attraction stabilizes the ionic crystal.

For ionic crystal, we define radius ratio as
\begin{equation*}
    \text{Radius Ratio}=\frac{r+ }{r_-}
\end{equation*}
where $r_+$ is the radius of the cation and $r_-$ is the radius of the anion.
This ratio determines how many anions can stably surround a cation (or vice versa) without causing excessive overlap or leaving excessive empty space.
\begin{enumerate}
    \item $\text{RR}<0.155$. The cation is too small to fit stably in any interstitial--small holes between the usual atoms into which smaller atoms may be placed--void and the structure tends to be unstable or distorted.
    \item $0.155<\text{RR}<0.255$. The cation can just fit into a triangular planar site, producing a coordination number of 3. This is rare in three-dimensional crystals.
    \item $0.255<\text{RR}<0.414$. The cation can fit into a tetrahedral site, giving a coordination number of 4. A common example is the zinc blende $\ce{ZnS}$ structure.
    \item $0.414<\text{RR}<0.732$. The cation can fit into an octahedral site, giving a coordination number of 6. The canonical example is the sodium chloride $\ce{NaCl}$ structure.
    \item $0.732<\text{RR}<1$. The cation can fit into a cubic site, giving a coordination number of 8. The archetypal structure here is cesium chloride $\ce{CsCl}$.
\end{enumerate}
\begin{figure*}[]
    \centering
    \dfig{../../../Rss/QM/ManyBody/RRRepresentation}
    \caption*{Figure: Representation of the location of the interstitial. Respectively, linear, center of a triangle, center of a tetrahedral, center of an octahedral, and center of a cube.}
\end{figure*}

\subsubsection{Sodium chloride.}
The crystalline structure of sodium chloride arises from the electrostatic attraction between sodium $\ce{Na+}$ cations and chloride anions $\ce{Cl-}$.
The radius ratio is
\begin{equation*}
    \frac{r_{\ce{Na+ }}}{r_{\ce{Cl-}}}=\frac{0.097 \text{ nm}}{0.181 \text{ nm}}=0.536
\end{equation*}
The resulting solid adopts a face-centered cubic lattice with 6 coordination number.
The FCC structure, with $\ce{Cl-}$ ions at FCC positions and $\ce{Na+}$ at the four octahedral sites.
We can also consider this structure to be FCC with two ions--one $\ce{Na+}$ and one $\ce{Cl-}$--associated with each lattice point.
\begin{figure*}[]
    \centering
    \normfig{../../../Rss/QM/ManyBody/NaCl}
    \caption*{Figure: The sodium chloride structure, an FCC unit cell with two ions ($\ce{Na+}$ and $\ce{Cl-}$) per lattice point}
\end{figure*}

\subsubsection{Cesium chloride.}
Cesium chloride forms a BCC-like structure, but it is not a true BCC lattice.
Instead, it is an SC lattice of $\ce{Cs+}$ ions, with $\ce{Cl-}$ ions occupying the cube centers.
Recall that in BCC all atoms are identical.

The radius ratio is
\begin{equation*}
    \frac{r_{\ce{Zn^{+2} }}}{r_{\ce{S^{-2}}}}=0.732
\end{equation*}
and thus having coordination number of eight.

\subsubsection{Zinc Blende.}
Also called $\ce{ZnS}$.
Zinc blende is constructed by cation $\ce{Zn^{+2}}$ and anion $\ce{S^{-2}}$ each with radius ratio
\begin{equation*}
    \frac{r_{\ce{Zn^{+2} }}}{r_{\ce{S^{-2}}}}=\frac{0.097 \text{ nm}}{0.181 \text{ nm}}=0.536
\end{equation*}
which yield coordination number of four.
The FCC structure, with $\ce{S}$ anions at the normal lattice points and Zn cations at half of the tetrahedral sites, can accommodate the restrictions of both charge balance and coordination number.
A variety of materials, including the semiconductor $\ce{GaAs}$ and many other III-V semiconductors have this structure.
\begin{figure*}[]
    \centering
    \normfigL{../../../Rss/QM/ManyBody/ZincBlende}
    \caption*{Figure: The zinc blende unit cell and the plan view}
\end{figure*}

\subsection{Bragg Theory}
X-ray diffraction is used to determine the structure of a crystal.
One of the analyzation is described by Bragg is his law.
Bragg treated the problem as the reflection of the incoming X-rays at flat crystal planes.
A collimated beam of monochromatic X-rays hits the crystal.
Recall that the interference condition is $n\lambda$. From the figure, $AB=d\sin \theta$.
As such,
\begin{equation*}
    n\lambda=2d\sin\theta
\end{equation*}
\begin{figure*}[]
    \centering
    \normfigL{../../../Rss/QM/ManyBody/Bragg}
    \caption*{Figure: Construction for the derivation of the Bragg condition.}
\end{figure*}

\subsection{Fourier Analysis}
A crystal is a periodic arrangement of atoms in real space, which means any property that repeats with the lattice, such as electron density $\rho(\mathbf{r})$ can be expressed as a Fourier series over reciprocal lattice vectors:
\begin{equation*}
    \rho(\mathbf{r}) = \sum_{\mathbf{G}} \rho_{\mathbf{G}} \, e^{i \mathbf{G} \cdot \mathbf{r}}
\end{equation*}
where the sum runs over all reciprocal lattice vectors $\mathbf{G} = h \mathbf{b}_1 + k \mathbf{b}_2 + l \mathbf{b}_3$.
The Fourier coefficients are
\begin{equation*}
    \rho_{\mathbf{G}} = \frac{1}{V} \int_{\text{unit cell}} \rho(\mathbf{r}) \, e^{-i \mathbf{G} \cdot \mathbf{r}} \, d^3 r.
\end{equation*}
Here, $V$ is the volume of the unit cell.
The reciprocal lattice vectors $\mathbf{G}$ form the natural basis for the Fourier decomposition of any lattice-periodic function, which is why diffraction experiments directly probe the reciprocal lattice.

The key point is that the periodicity of the crystal in real space translates to discrete points in reciprocal space, which simplifies the analysis of diffraction and electronic structure.

\subsubsection{1D case.}
Consider a one-dimensional lattice with spacing $a$, with lattice points at
\begin{equation*}
    x_n = n a
\end{equation*}
A reciprocal lattice vector $G_m$ may be determined with
\begin{align*}
    e^{i G_m x_n} & =  1 \\
    e^{i G_m n a} & =  1
\end{align*}
This condition holds for all integers $n$ if the exponent is an integer multiple of $2\pi$.
Thus,
\begin{equation*}
    i G_m na=2\pi m \implies G_m = \frac{2 \pi m}{a}
\end{equation*}
with integer $m$.
A periodic function on the lattice, such as the electron density $\rho(x)$, can then be expressed as a Fourier series over reciprocal lattice vectors:
\begin{equation*}
    \rho(x) = \sum_{m=-\infty}^{\infty} \rho_m \, e^{i G_m x}, \quad
    \rho_m = \frac{1}{a} \int_0^a \rho(x) \, e^{-i G_m x} \, dx
\end{equation*}
\end{document}